\documentclass[11pt]{report}

% --- Basics ---
\usepackage[a4paper,margin=1in]{geometry}
\usepackage[hidelinks]{hyperref}
\usepackage{graphicx}
\usepackage{float}      % for [H] figure placement when needed
\usepackage{caption}
\usepackage{booktabs}   % nice tables (optional)

% Optional: if you want better crossrefs like "Figure 3.2"
% \usepackage[nameinlink,noabbrev]{cleveref}

% Optional: captions for many screenshots can get repetitive
% \captionsetup[figure]{font=small}

% --- Screenshot helpers (put in the preamble) ---
\usepackage{graphicx}
\usepackage{float}      % enables [H]
\usepackage{caption}

% Optional: common folder for images
\newcommand{\imgdir}{figs/}

% Main screenshot macro:
%  - #1: filename (relative to \imgdir unless you include a path)
%  - #2: caption text
%  - #3: label (without "fig:"; the macro adds it)
%  - #4: width (optional; default 0.98\linewidth)
\newcommand{\screenshot}[4][0.98\linewidth]{%
    \begin{figure}[H]
        \centering
        \includegraphics[width=#1]{\imgdir#2}
        \caption{#3}
        \label{fig:#4}
    \end{figure}
}

% Variant: screenshot with a short caption for LoF + long caption for the page
%  - #1 width (optional; default 0.98\linewidth)
%  - #2 filename
%  - #3 short caption (List of Figures)
%  - #4 long caption (printed under figure)
%  - #5 label key
\newcommand{\screenshotLF}[5][0.98\linewidth]{%
    \begin{figure}[H]
        \centering
        \includegraphics[width=#1]{\imgdir#2}
        \caption[#3]{#4}
        \label{fig:#5}
    \end{figure}
}

% Quick reference helper for figures
\newcommand{\figref}[1]{Figure~\ref{fig:#1}}


\title{Transgraphier\\{\large User Manual}}
\author{Fernando Cacciola}
\date{Version 2.3 --- \today}

\begin{document}
    \maketitle
    \tableofcontents
    
    % -------------------------
    \chapter{Introduction}
    \section{What the App Does}
       
       Given an \textbf{EVP} audio file that is \textit{specifically} is expected
       to contain a text message \textit{encoded in binary} (like computers do),
       but using \textbf{Tap Code} as encoding system, applies a sequence of audio processing steps in order to attempt to extract a Text Message from it.
       
    \section{Typical Workflow at a Glance}
    
\begin{itemize}
    \item Open an audio file containing an \textbf{EVP}.
         
    \item Click on Process (or Re-Process if this one file has already been processed)
         
    \item Look at the \textbf{TOP-RIGHT} panel.\\
          If will either say that ``\textit{no message was decoded}'' OR\\
          It will contain the text message that was decoded from the \textbf{EVP} audio.
         
    \item The TOP-LEFT panel will show a general log detailing what\\
          the processing was doing.
         
    \item You can also look at the sequence of ``Filter Outputs'', right below\\
           the Input Wave Form.\\
           Details about the result of each processing step are shown there.
\end{itemize}
    

\end{document}
